% =========================================================
\section{Introduction}
\label{sec:introduction}
% =========================================================

Gradient descent with a finite step size is naturally a discrete
dynamical system,
\begin{equation}
x_{n+1}
=
x_n-\eta\nabla V(x_n),
\label{eq:intro-GD}
\end{equation}
and its behavior near the boundary of linear stability can differ
substantially from that suggested by the continuous gradient-flow
approximation; see, for example,
\cite{RobbinsMonro1951,Polyak1987,Nesterov2018,BottouCurtisNocedal2018}
for the classical optimization background.  In particular, when a nondegenerate local minimum
$x_*$ has maximal Hessian eigenvalue
\[
\lambda_*=
\lambda_{\max}(D^2V(x_*)),
\]
the classical local stability threshold is
\begin{equation}
\eta_f
=
\frac{2}{\lambda_*}.
\label{eq:intro-flip-threshold}
\end{equation}
At this value, the multiplier corresponding to the sharpest
eigendirection reaches $-1$, and, under the usual nondegeneracy
conditions, a flip bifurcation may occur.
This is the classical simple flip-bifurcation mechanism for smooth
maps; see, for example,
\cite{Kuznetsov2004,GuckenheimerHolmes1983,HaleKocak1991}.
The dynamics of gradient descent near this threshold have recently
received considerable attention, partly motivated by the
edge-of-stability phenomenon observed in large-step
neural-network optimization
\cite{CohenEtAl2021}.  Beyond
linear stability, higher-order derivatives of the objective determine
whether the loss of stability gives rise to stable or unstable
oscillations.  In the multivariate setting, the nonlinear coefficient
governing the local period-two branch depends not only on derivatives
along the maximal-curvature eigendirection, but also on its coupling to
the remaining Hessian eigendirections; see
\cite{MulayoffStich2026,Gan2026,Litman2026}.
These developments provide a
precise smooth bifurcation description of oscillatory gradient-descent
dynamics near a minimum.

Gradient clipping is a standard mechanism for controlling large
updates in optimization and has been studied from the viewpoints of
training stability, adaptivity, and convergence
\cite{PascanuMikolovBengio2013,ZhangEtAl2020,KoloskovaHendrikxStich2023,Goodfellow2016}.
The situation changes qualitatively when the gradient is clipped.
We consider the radial clipped gradient map
\begin{equation}
F_{\eta,c}(x)
=
x-\eta\mathcal C_c^\diamond(\nabla V(x)),
\label{eq:intro-clipped-map}
\end{equation}
where
\begin{equation}
\mathcal C_c^\diamond(g)
=
\min
\left\{
1,\frac{c}{\|g\|_\diamond}
\right\}g,
\qquad
c>0.
\label{eq:intro-clipping}
\end{equation}
Here $\|\cdot\|_\diamond$ is a smooth norm away from the origin.
The standard Euclidean clipping rule is recovered by taking
$\|\cdot\|_\diamond=\|\cdot\|_2$.

The clipping rule introduces the state-dependent switching
hypersurface
\begin{equation}
\Sigma_c
=
\left\{
x:
\|\nabla V(x)\|_\diamond=c
\right\}.
\label{eq:intro-switching-surface}
\end{equation}
Inside
\[
\Omega_c^-
=
\{x:\|\nabla V(x)\|_\diamond<c\},
\]
the map \eqref{eq:intro-clipped-map} coincides exactly with ordinary
gradient descent, whereas in
\[
\Omega_c^+
=
\{x:\|\nabla V(x)\|_\diamond>c\}
\]
the step magnitude is saturated.  Thus clipping converts a smooth
gradient map into a continuous piecewise-smooth dynamical system whose
switching boundary is not imposed externally: it is generated
endogenously by the same potential $V$ that generates the smooth
dynamics.  Piecewise-smooth maps and their border-collision
bifurcations have been studied extensively; see
\cite{Filippov1988,NusseYorke1992,diBernardo2008,diBernardoEtAl2008,
diBernardoHogan2010,Simpson2010,Simpson2016,LeineNijmeijer2004}.

This observation leads to a basic separation of mechanisms.  Since
\[
\nabla V(x_*)=0,
\]
for every fixed $c>0$ there exists a neighborhood of $x_*$ entirely
contained in $\Omega_c^-$.  Consequently, clipping cannot change the
local germ of the map at the minimum.  In particular, the fixed point,
its linearization, its first loss-of-stability threshold, the
center-manifold reduction, and the smooth flip coefficient are all
identical to those of ordinary gradient descent.  The first genuinely
clipping-induced effect must therefore occur at finite amplitude,
after the bifurcating orbit has grown far enough to reach
$\Sigma_c$.

The aim of this paper is to characterize this first
smooth-to-switching transition.

% ---------------------------------------------------------
\subsection{Main question}
\label{subsec:intro-main-question}
% ---------------------------------------------------------

Assume that $x_*$ is a nondegenerate strict local minimum,
that the maximal Hessian eigenvalue $\lambda_*$ is simple, and that
the ordinary gradient map undergoes a supercritical flip bifurcation
at
\[
\eta_f=\frac{2}{\lambda_*}.
\]
For
\[
\mu=\eta-\eta_f>0
\]
small, the smooth map possesses a stable period-two cycle
\begin{equation}
\mathcal O_2(\mu)
=
\{x_-(\mu),x_+(\mu)\}.
\label{eq:intro-cycle}
\end{equation}

The central question is:

\begin{quote}
At what parameter value does the smooth period-two branch first reach
the clipping boundary, and what dynamical structure is created when
this contact occurs?
\end{quote}

For a generic flip bifurcation, the oscillation amplitude grows on the
square-root scale
\[
\|x_\pm(\mu)-x_*\|
\asymp
\sqrt{\mu},
\]
whereas the local switching boundary generated by
\[
\|\nabla V(x)\|_\diamond=c
\]
lies at distance $O(c)$ from the critical point.  This suggests the
matching law
\[
\mu_{\mathrm{sc}}
\asymp
c^2.
\]
However, the main issue is not merely the exponent.  General
period-doubling--border-collision interactions are already known to
produce tangencies between smooth bifurcation and switching loci.
What is special here is the gradient-generated nature of both the
smooth map and the switching surface.

This additional structure allows us to derive the separation
coefficient explicitly from the local derivatives of $V$, to prove
that the two points of the period-two orbit reach the switching
surface simultaneously, and to determine a universal spectral
degeneracy in the fully clipped regime.

% ---------------------------------------------------------
\subsection{Main results}
\label{subsec:intro-results}
% ---------------------------------------------------------

Let
\[
H_*
=
D^2V(x_*),
\qquad
T_*
=
D^3V(x_*),
\qquad
Q_*
=
D^4V(x_*),
\]
and let
\[
H_*v_*=\lambda_*v_*,
\qquad
\|v_*\|_2=1.
\]
Define
\begin{equation}
g_*
=
T_*[v_*,v_*,\cdot]^\sharp
\label{eq:intro-g}
\end{equation}
and
\begin{equation}
\Gamma_*
=
3T_*[v_*,v_*,H_*^{-1}g_*]
-
Q_*[v_*,v_*,v_*,v_*].
\label{eq:intro-Gamma}
\end{equation}
The condition
\[
\Gamma_*>0
\]
is the supercritical flip condition used throughout the paper.

Our first result is a local identity principle.  For every fixed
$c>0$, the clipped and unclipped maps coincide on a neighborhood of
$x_*$.  Hence
\begin{equation}
DF_{\eta,c}(x_*)
=
I-\eta H_*,
\label{eq:intro-local-linearization}
\end{equation}
and the first local stability threshold is exactly
\begin{equation}
\boxed{
\eta_f
=
\frac{2}{\lambda_*},
}
\label{eq:intro-threshold-box}
\end{equation}
independently of $c$.  More strongly, the finite jets of the two maps
at $x_*$ agree, so the local center dynamics and smooth flip
coefficient are also clipping-independent.

Second, we derive the leading amplitude of the smooth period-two
branch.  For
\[
\mu=\eta-\eta_f\to0^+,
\]
one has
\begin{equation}
x_\pm(\mu)
=
x_*
\pm
A_*\sqrt{\mu}\,v_*
+
O(\mu),
\label{eq:intro-amplitude}
\end{equation}
where
\begin{equation}
\boxed{
A_*^2
=
\frac{3\lambda_*^2}{\Gamma_*}.
}
\label{eq:intro-A}
\end{equation}

Third, the gradient structure gives an exact identity that is stronger
than the local normal-form expansion.  For every genuine period-two
cycle of ordinary gradient descent,
\begin{equation}
\boxed{
\nabla V(x_+)
=
-\nabla V(x_-).
}
\label{eq:intro-gradient-balance}
\end{equation}
Thus the two points have exactly the same clipping exposure for every
norm:
\[
\|\nabla V(x_+)\|_\diamond
=
\|\nabla V(x_-)\|_\diamond.
\]
Consequently, the first interaction with a radial clipping boundary is
necessarily a simultaneous two-point contact.

Along the bifurcating branch, the common exposure satisfies
\begin{equation}
\|\nabla V(x_\pm(\mu))\|_\diamond
=
B_*\sqrt{\mu}
+
O(\mu^{3/2}),
\label{eq:intro-exposure}
\end{equation}
where
\begin{equation}
B_*
=
\lambda_*A_*
\|v_*\|_\diamond.
\label{eq:intro-B}
\end{equation}
Inverting the contact equation
\[
\|\nabla V(x_\pm(\mu_{\mathrm{sc}}))\|_\diamond=c
\]
gives our principal separation theorem:
\begin{equation}
\boxed{
\eta_{\mathrm{sc}}(c)-\eta_f
=
K_\diamond c^2
+
O(c^4),
\qquad
c\to0^+,
}
\label{eq:intro-main-law}
\end{equation}
with the explicit coefficient
\begin{equation}
\boxed{
K_\diamond
=
\frac{\Gamma_*}{
3\lambda_*^4
\|v_*\|_\diamond^2
}.
}
\label{eq:intro-K}
\end{equation}

Equivalently, if
\[
\{v_i\}_{i=1}^d
\]
is an orthonormal Hessian eigenbasis with eigenvalues
$\lambda_i$, then
\begin{equation}
K_\diamond
=
\frac{
3\displaystyle\sum_{i=1}^d
\frac{
T_*[v_*,v_*,v_i]^2
}{
\lambda_i
}
-
Q_*[v_*^4]
}{
3\lambda_*^4
\|v_*\|_\diamond^2
}.
\label{eq:intro-K-eigen}
\end{equation}

Thus the same high-order local quantity that controls the birth and
stability of the smooth oscillation also determines the leading
distance in parameter space to the first clipping event.

The radial structure imposes further restrictions after contact.
A period-two orbit of the clipped map cannot have one point in
$\Omega_c^-$ and the other in $\Omega_c^+$.  Hence a mixed
smooth--clipped period-two itinerary is impossible.  If a period-two
orbit persists beyond contact, both points must lie in the clipped
region and satisfy the exact amplitude constraint
\begin{equation}
\boxed{
\|y_+-y_-\|_\diamond
=
\eta c.
}
\label{eq:intro-clipped-amplitude}
\end{equation}

Finally, every fully clipped period-two orbit is necessarily
nonhyperbolic.  If
\[
M_c
=
DF_{\eta,c}(y_-)
DF_{\eta,c}(y_+)
\]
denotes its monodromy matrix, then
\begin{equation}
\boxed{
1\in\sigma(M_c).
}
\label{eq:intro-unit-multiplier}
\end{equation}
This unit multiplier follows from the fact that radial clipping
annihilates infinitesimal changes in gradient magnitude while
retaining only tangent variations of the gradient direction.

A completely solvable planar gradient model shows that this
degeneracy can have a genuine nonlinear manifestation: after the
simultaneous switching contact, the isolated attracting smooth
period-two cycle is replaced locally by a one-parameter family of
fully clipped period-two cycles, transversely attracting but neutral
along the family.

\subsection{Relation to existing bifurcation theory}

The present analysis lies at the intersection of two established
lines of work.

The first is smooth local bifurcation theory for maps.
Near a simple multiplier $-1$, center-manifold and normal-form
reductions yield the classical flip bifurcation and the
square-root scaling of the emerging period-two branch; see,
for example,
\cite{Kuznetsov2004,GuckenheimerHolmes1983,ChowHale1982,
Wiggins2003,Arnold1983,Sotomayor1973,HaleKocak1991,PalisDeMelo1982}.

More recently, nonlinear oscillatory dynamics of gradient descent
near the edge of stability have been studied directly.
Mulayoff and Stich \cite{MulayoffStich2026} derived a multivariate
stability criterion for nonlinear gradient-descent oscillations
near a minimum in terms of higher-order derivatives of the
objective.
Gan \cite{Gan2026} formulated edge-of-stability dynamics in a
bifurcation-theoretic framework, in which the relevant normal
dynamics undergo a flip bifurcation governed by the corresponding
first Lyapunov coefficient.
Relatedly, Litman \cite{Litman2026} characterized fixed points and
period-two orbits of gradient descent through an exact two-step
variational formulation and obtained a local reduction in terms
of the half-amplitude of the oscillation.

Some ingredients of the smooth analysis developed below therefore
belong to standard flip-bifurcation theory or have close analogues
in recent analyses of nonlinear gradient-descent dynamics.
We nevertheless include self-contained derivations, both to fix
notation and because the resulting coefficients enter explicitly
into the subsequent switching analysis.

The second relevant line of work is the theory of piecewise-smooth
maps and border-collision bifurcations
\cite{NusseYorke1995,BanerjeeGrebogi1999,Feigin1970,diBernardo2008,
Simpson2010,Simpson2016}.
In particular, Simpson and Meiss
\cite{SimpsonMeiss2008,SimpsonMeiss2009,SimpsonMeiss2009Nonlinearity,
SimpsonMeiss2010,SimpsonMeiss2012} studied the codimension-two
interaction between period-doubling and border-collision
bifurcations in general piecewise-smooth continuous maps. They
showed that the border-collision locus of the period-doubled
solution can be tangent to the period-doubling locus near the
codimension-two point. Thus quadratic separation between smooth
bifurcation and switching loci is already a feature of generic
period-doubling--border-collision interactions.

Our contribution is not the generic existence of such a tangency,
nor the smooth flip mechanism itself. Rather, we exploit the
additional variational structure of clipped gradient maps.
The switching surface
\[
\Sigma_c=\{x:\|\nabla V(x)\|_\diamond=c\}
\]
is generated endogenously by the same potential that determines
the smooth gradient dynamics. Moreover, the period-two equations
impose exact gradient-balance identities.

These additional constraints yield structural consequences that
are absent from a generic externally switched map. In particular,
we obtain:
\begin{enumerate}
    \item clipping-independent local flip dynamics;
    \item an explicit tensor formula for the flip-to-switching
    separation coefficient;
    \item simultaneous switching contact of the two points of a
    smooth period-two orbit;
    \item nonexistence of mixed smooth--clipped period-two
    itineraries;
    \item an exact unit Floquet multiplier for every fully clipped
    period-two orbit.
\end{enumerate}

It is these gradient-induced rigidity properties, rather than the
standard smooth flip bifurcation or the generic tangency mechanism,
that form the main focus of the present work.
% ---------------------------------------------------------
\subsection{Organization of the paper}
\label{subsec:intro-organization}
% ---------------------------------------------------------

The paper is organized as follows.

Section~\ref{sec:setting} introduces the clipped gradient map,
the switching hypersurface, and the standing assumptions.

Section~\ref{sec:local-structure} proves the local identity between
clipped and ordinary gradient descent and establishes the
clipping-independent stability threshold.

Section~\ref{sec:flip-bifurcation} derives the supercritical
period-two branch and its explicit leading amplitude.

Section~\ref{sec:switching-geometry} develops the exact gradient
balance and switching geometry along the period-two branch.

Section~\ref{sec:separation-law} proves the quadratic
flip-to-switching separation law and derives its explicit coefficient.

Section~\ref{sec:post-contact} studies the fully clipped regime and
proves the universal unit-multiplier theorem.

Section~\ref{sec:planar-model} presents a completely solvable planar
model that realizes the full transition analytically.
